\begin{textsample}{NEG Dim 7 – global\_south\_2024\_03 – Score -4.00 – global\_south\_2024\_03/106924665.txt}  \label{ex:f7_neg_423}
We should listen to the Global South ' s outcry on the Gaza war An urgent solution is necessary, although there seems to be little hope of one anytime soon, despite ongoing talks in Qatar aimed at freeing the hostages and pausing fighting for several weeks. ( AFP ) Summary Calls for a ceasefire and increase in aid need to be taken up urgently.
It ' s emerging as the world ' s conscience in a war that ' s gone too far.
There is a rising chorus of voices, mostly from the Global South, calling for an immediate ceasefire in Gaza and the urgent provision of aid to a population that is facing imminent famine.
Is \textbf{anybody} listening?
Singapore ' s foreign minister Vivian Balakrishnan is the latest to intervene, telling Israel that it had"gone too far"in its war against Hamas in Gaza."I have communicated that both to the prime minister, to the foreign minister, and to the other Israelis whom we have met,"he said last week on Wednesday.
The tiny island state @ @ @ @ @ @ @ @ @ @ this issue -- just under a fifth of the population is Muslim.
Like other leaders, Balakrishnan condemned the \textbf{horrific} events of 7 October, when Hamas militants killed 1,200 people and took more than 200 hostage.
Since then, Israel ' s military has killed more than 31,000 Palestinians, the majority of them women and children, the health ministry in the Hamas-run territory says.
An urgent solution is necessary, although there seems to be little hope of one anytime soon, despite ongoing talks in Qatar aimed at freeing the hostages and pausing fighting for several weeks.
A famine will escalate the scale of death exponentially.
Which is why the Global South -- a collection of post-colonial and developing countries that by some estimates represent 88 \% of the world ' s population -- has taken on the cause of Gaza so passionately : It feels like no one is listening.
Australia and Southeast Asian nations have also called for a lasting ceasefire, and the need to urgently upscale aid deliveries.
Brazilian President Luiz Inacio Lula da Silva has advocated @ @ @ @ @ @ @ @ @ @ United Nations.
There are domestic issues at play -- Brazil is home to an estimated 60,000 Palestinian immigrants and refugees, including their descendants.
Lula has described Israel ' s prosecution of the war as ' genocide ' -- a position taken by South Africa, too, in its case in the International Court of Justice, but also, further afield, by Ireland.
At a St.
Patrick ' s Day celebration at the White House this week, Leo Varadkar -- an outspoken critic of the impact on civilians of Israel ' s war -- who until last week was Ireland ' s prime minister, called for a ceasefire and an increase in humanitarian aid.
This conflict has particular resonance for the Irish, given their own history of resistance to British rule.
Ireland is one of the more supportive European nations to Palestinians and was the first to call for a Palestinian state in 1980."When I travel the world, leaders often ask me why the Irish people have such empathy for the Palestinian people,"Varadkar said. @ @ @ @ @ @ @ @ @ @ in their eyes."These comments have shone a light on the depth of feeling about the inhumanity of Israel ' s war in Gaza, not just among the political elite, but also with the youthful populations of the US and the UK including Scotland and Europe.
Their words bear weight and are worth listening to.
US President Joe Biden, one of Israel ' s staunchest supporters and the global leader most able to influence the course of the ongoing war -- by restricting weapons sales and pushing for an immediate ceasefire in the United Nations -- should pay attention to this growing tide of dissent.
Perhaps no one has said this as eloquently as Malaysia ' s Prime Minister Anwar Ibrahim, when he asked :"Where have we \textbf{thrown} our humanity, why this hypocrisy?"What he said struck a chord.
In an interview on 16 March with DW News, the German public broadcaster, he asked why there is"selective amnesia"-- suggesting that the difference in treatment and attitude from some Western governments @ @ @ @ @ @ @ @ @ @ hospitals, universities, mosques and \textbf{churches} in Gaza, compared to Ukraine, has something to do with skin colour and the Islamic faith.
Of course, Anwar Ibrahim has ideological reasons to support the plight of Palestinians, and has refused to cut ties with Hamas ' s political wing.
Popular opinion is also causing him to speak out, given the strength of support domestically.
It would be easy to dismiss his demands given his roots in Islamic student politics and the dependence of his government on the goodwill of the strong Parti Islam Se-Malaysia ( PAS ) party.
But he has long been an articulate voice on human rights, and what he is asking for is a transparent, consistent and cohesive voice on these conflicts.

% matched lemmas: anybody, church, horrific, throw
\end{textsample}
